\section{Qué parte del resto se estima}
Ponemos $K=\lceil\sqrt n\rceil$, $j=\lceil tK\rceil\le M=\lceil\delta n\rceil$,
$s=1/j$ y $\E_* Z=\E_{\bar\nu}[Z;G_j,S_M]$. Suponemos las entradas uniformes
de difusión, brecha, estabilidad de menores y desarrollo de torres de R8--R9
usadas en la entrega de las cuatro marcas, y la prescripción conjunta de llegada
$\Pr_*(h\text{ coordenadas retenidas})\le C_h(j/n)^h$ para $h$ fijo.
Los resultados siguientes son deducciones bajo estas hipótesis.

En una misma máscara $A$ de torres, para $Y$ subconjunto de las casillas se activa
la capacidad uno en cada diagonal y antidiagonal sólo sobre sus sitios de $Y$.
Las casillas de $A\setminus Y$ siguen disponibles para las torres. Sea
$e(Y)=\log(Q_Y/\per A)-(L_Y-L_T)$. Para un conjunto de $v$ sitios $S$ escribimos
\[
 [S]e=\sum_{Y\subseteq S}(-1)^{v-|Y|} e(Y).
\]
En cada comparación se conserva la máscara y se reoptimiza el predictor.
Para cada $v$ fijo y $j$ grande existen todas las evaluaciones de tamaño $\le v$:
la probabilidad de un conflicto entre $v$ marcas es $O_v(j^{-2})$ por R2;
el punto pequeño es la continuación que se construye en el apartado 3.
No se presupone la existencia de predictores para toda subfamilia de tamaño creciente.

A cada sitio corresponde una arista entre su diagonal $D$ y su antidiagonal $A'$.
Dos sitios quedan conectados si se puede ir de uno al otro pasando por líneas
con una casilla marcada en común. Esta conectividad es geométrica. Los sitios de
dos componentes distintas pueden interactuar por la medida de torres.
\paragraph{Resultado.}
Para $v$ fijo, la suma absoluta en media de $[S]e$ sobre soportes diagonalmente
conexos de $v\ge2$ sitios es $O_v(1/n)$, uniformemente en los suelos de la banda.
La misma cota se aplica al resto de la entrega anterior después de descontar los
pesos de líneas individuales y de dos conflictos. Al restringir a soportes que usan
al menos dos líneas no unitarias, una línea individual tiene proyección cero.
Un peso de dos conflictos tiene proyección sobre $S$ sólo si su unión es exactamente
$S$: sobre $Y$ pesa su valor físico multiplicado por el indicador de esa unión en $Y$.
Por tanto, para $v\ge5$ desaparece y para $v\le4$ basta su cota absoluta anterior.
En el soporte enteramente colineal el resto es cero. Un sitio aislado anula la
proyección. Con cuatro sitios desconectados sin aislados sólo pueden quedar dos pares
y su resto es cero por definición, aunque sus coordenadas RC se solapen.
Así se cubre toda la parte de cuatro sitios del resto.
